\thispagestyle{empty}
\vspace*{\fill}
\begin{minipage}{128mm}
\fontsize{9}{13}\selectfont\color{otMuted}
{\otSemi\color{otInk}Physique-Chimie · Terminale spécialité · Livre complet}\par
\vspace{3mm}
Conception et mise en page : collection OptiTAB.\par
\otNbExosTotal\ exercices originaux et \otNbSujetsLettres\ sujets type bac avec leurs corrigés.\par
Charte graphique commune aux livres OptiTAB de mathématiques.\par
\vspace{5mm}
Édition du 21 octobre 2026.\par
Constitution et transformations de la matière, mouvement et interactions,
énergie, ondes et signaux, mesures et incertitudes.\par
\vspace{5mm}
Référence pédagogique : programme de l'enseignement de spécialité de physique-chimie
de la classe terminale de la voie générale, publié au Bulletin officiel spécial n° 8
du 25 juillet 2019 (NOR MENE1921249A), en vigueur depuis la rentrée 2020.
L'ordre des chapitres est un choix pédagogique de ce livre.\par
\vspace{2mm}
Épreuve de référence : note de service du 11 septembre 2026 (NOR MENE2622644N),
publiée au Bulletin officiel spécial n° 4 du 17 septembre 2026, applicable
à compter de la session 2027 du baccalauréat.\par
\vspace{5mm}
Les schémas sont construits à partir des données des énoncés ; ceux qui ne
respectent pas d'échelle sont signalés. Les données réelles sont sourcées dans
les exercices concernés ; les autres situations sont construites pour l'entraînement.
Les valeurs numériques sont arrondies avec un nombre de chiffres significatifs adapté aux données.\par
\vspace{5mm}
© OptiTAB, 2026.\par
\otSite
\end{minipage}
\clearpage
